\par\Needspace{9\baselineskip}\bigskip
\pdfbookmark[3]{Darstellung, [23]}{pallas-II-article23}
\begin{center}[23.]\end{center}
\begingroup
\setlength{\lineskip}{4pt}\setlength{\lineskiplimit}{4pt}
\noindent\hspace*{\parindent}Mit dieser Methode bestimmt man die Elemente \(\AscendingNode+\delta\AscendingNode,\ i+\delta i,\ \varphi+\delta\varphi,\ \tilde\omega+\delta\tilde\omega,\ N+\delta N,\ n+\delta n\). In der Entwicklung der momentanen Änderung von \(n+\delta n\) zeigt sich, dass die Koeffizienten aller Glieder mit \(i=0\) gleich null werden. Daraus folgt, dass weder die mittlere Bewegung noch die große Achse eine eigentliche säkulare Änderung enthält, zumindest wenn man sich auf die erste Potenz der störenden Kraft beschränkt. Diese Elemente enthalten auch keine von \(L\) unabhängigen periodischen Glieder. Um dies unmittelbar zu beweisen, greifen wir die Formel \(m\mathrm{d}\dfrac{1}{a}\bigl(\text{oder vielmehr }m\mathrm{d}\dfrac{1}{a+\delta a}\bigr)=2(\xi\mathrm{d}x+\eta\mathrm{d}y+\zeta\mathrm{d}z)\) wieder auf. Man sieht leicht, dass, wenn man
\[
m\left\{\frac{xx'+yy'+zz'}{r'^{8}}-\frac{1}{\rho}\right\}=\Omega
\]
\SourceLeaf{473}setzt, die Kräfte \(\xi,\eta,\zeta\) durch die partiellen Differentiale ausgedrückt werden
\[
\xi=\mu\left(\frac{\mathrm{d}\Omega}{\mathrm{d}x}\right)\mathpunct{\raisebox{.45ex}{,}}\qquad
\eta=\mu\left(\frac{\mathrm{d}\Omega}{\mathrm{d}y}\right)\mathpunct{\raisebox{.45ex}{,}}\qquad
\zeta=\mu\left(\frac{\mathrm{d}\Omega}{\mathrm{d}z}\right)\mathpunct{\cdot}
\]
Somit ist \(\mathrm{d}\dfrac{2\mu}{a+\delta a}\) das Differential von \(m\Omega\), wenn man die auf den störenden Planeten bezogenen Größen als konstant betrachtet. Nun ist klar, dass \(\Omega\) eine \textit{periodische} Funktion von \(\boldsymbol{L}\) und \(\boldsymbol{L}'\) ist. Entwickelt man sie also in Glieder der Form \(k\cos(iL+i'L')\) und \(k\sin(iL+i'L')\), so enthält ihr Differential, bei konstant gehaltenem \(L'\) gebildet, kein von \(L\) unabhängiges Glied. Da außerdem gilt
\[
\begin{aligned}
n+\delta n&=\BareRoot{m}.(a+\delta a)^{-\frac{3}{2}}\\
\mathrm{d}(n+\delta n)&=-\tfrac{3}{2}\BareRoot{m}.(a+\delta a)^{-\frac{5}{2}}\mathrm{d}(a+\delta a)=\tfrac{3}{2}\BareRoot{\dfrac{m}{a+\delta a}}\cdot\mathrm{d}\frac{1}{a+\delta a}\mathpunct{\raisebox{.45ex}{,}}
\end{aligned}
\]
erhält man bei Berücksichtigung allein der Glieder erster Ordnung
\[
\mathrm{d}(n+\delta n)=3\frac{\mu}{\BareRoot{ma}}\cdot\mathrm{d}\Omega,
\]
wobei das Differential \(\mathrm{d}\Omega\) in Bezug auf \(L\) gebildet wird. Daraus folgt, dass die Änderung \(\mathrm{d}(n+\delta n)\) die oben angegebene Form hat.
\par
Ist \(\bar n\) die durch Integration von \(\mathrm{d}(n+\delta n)\) eingeführte Konstante, so ist klar, dass das der mittleren Bewegung Entsprechende, das heißt das Integral \(\int n\mathrm{d}t\), sich zusammensetzt
\begin{enumerate}\setlength{\itemsep}{0pt}\setlength{\parsep}{0pt}
\item[1)] aus periodischen Gliedern
\item[2)] aus einem zur\raisebox{.45ex}{-} Zeit proportionalen Glied, \(=\bar n t\)
\item[3)] aus einem konstanten Teil.
\end{enumerate}
\par
Verbindet man also diese mittlere Bewegung mit der veränderlichen Epoche, die sich ihrerseits zusammensetzt
\begin{enumerate}\setlength{\itemsep}{0pt}\setlength{\parsep}{0pt}
\item[1)] aus periodischen Gliedern
\item[2)] aus einer säkularen Änderung, die wir mit \(\bar\varepsilon t\) bezeichnen
\item[3)] aus einem konstanten Teil,
\end{enumerate}
\noindent so hat die mittlere Länge die Form \(\bar L+(\bar n+\bar\varepsilon)t+P\), wobei \(P\) die periodischen Teile bezeichnet.
\par
Auf diese Weise bleiben sechs konstante Glieder übrig, die wir die \textit{mittleren Werte der Elemente} nennen. Verbindet man sie mit den \SourceLeaf{474}periodischen Teilen und den säkularen Änderungen, so erhält man die vollständigen Werte der Elemente für jeden Zeitpunkt.
\par
Für die konstanten Teile der Elemente \(\AscendingNode,\ i,\ \varphi,\ \tilde\omega\), nach denen man die Störungen berechnet, wählt man zweckmäßigerweise die mittleren Werte selbst, also die durch Integration eingeführten Konstanten, und für [diejenigen von] \(L\) und \(n\) die Größen \(\bar L\) und \(\bar n+\bar\varepsilon\). Man sieht jedoch, dass dafür in der Praxis zwei Rechnungen unerlässlich sind. Nachdem man die erste Rechnung mit beliebigen Näherungswerten der Elemente [ausgeführt] hat, bestimmt man nach bekannten Methoden die Werte, die man den sechs Konstanten geben muss, um die beste Übereinstimmung mit den Beobachtungen zu erhalten, und wiederholt die Rechnung mit diesen Werten. Sofern die ersten Werte nicht zu weit von den wahren entfernt waren, wird man die Rechnung gewiss kein drittes Mal ausführen müssen.
\par\endgroup
